\chapter{Vectors and Matrices}

\section{Euclidean Vectors and Linear Combinations}

A \textbf{vector} in $\R^n$ is an ordered list of $n$ real numbers. Geometrically, a vector represents a directed line segment possessing both magnitude and direction. Algebraically, it is represented as a column matrix.

\begin{definition}[Euclidean Vector Space $\R^n$]
    For a positive integer $n$, the set of all $n$-tuples of real numbers is called \textbf{$n$-dimensional Euclidean space}, denoted by $\R^n$:
    \[
        \R^n = \set{\vct{v} = \begin{pmatrix} v_1 \\ v_2 \\ \vdots \\ v_n \end{pmatrix}}{v_1, v_2, \dots, v_n \in \R}
    \]
    The real numbers $v_1, v_2, \dots, v_n$ are called the \textbf{components} (or coordinates) of the vector $\vct{v}$.
\end{definition}

\subsection*{Operations on Vectors}
The algebraic structure of $\R^n$ is established by two fundamental operations: \textbf{vector addition} and \textbf{scalar multiplication}.

\begin{definition}[Vector Addition]
    Let $\vct{u}, \vct{v} \in \R^n$ be vectors given by:
    \[
        \vct{u} = \begin{pmatrix} u_1 \\ u_2 \\ \vdots \\ u_n \end{pmatrix}, \quad 
        \vct{v} = \begin{pmatrix} v_1 \\ v_2 \\ \vdots \\ v_n \end{pmatrix}
    \]
    The \textbf{sum} $\vct{u} + \vct{v}$ is the vector in $\R^n$ obtained by adding corresponding components:
    \[
        \vct{u} + \vct{v} = \begin{pmatrix} u_1 + v_1 \\ u_2 + v_2 \\ \vdots \\ u_n + v_n \end{pmatrix}
    \]
\end{definition}

\begin{example}[Vector Addition]
    Let $\vct{u} = \begin{pmatrix} 2 \\ 1 \end{pmatrix}$ and $\vct{v} = \begin{pmatrix} 4 \\ 3 \end{pmatrix}$ in $\R^2$. Their sum is:
    \[
        \vct{u} + \vct{v} = \begin{pmatrix} 2 \\ 1 \end{pmatrix} + \begin{pmatrix} 4 \\ 3 \end{pmatrix} = \begin{pmatrix} 2 + 4 \\ 1 + 3 \end{pmatrix} = \begin{pmatrix} 6 \\ 4 \end{pmatrix}
    \]
\end{example}

\begin{definition}[Scalar Multiplication]
    Let $\vct{v} \in \R^n$ and let $c \in \R$ be a real number (called a \textbf{scalar}). The \textbf{scalar product} $c\vct{v}$ is obtained by multiplying each component of $\vct{v}$ by $c$:
    \[
        c\vct{v} = c\begin{pmatrix} v_1 \\ v_2 \\ \vdots \\ v_n \end{pmatrix} = \begin{pmatrix} c v_1 \\ c v_2 \\ \vdots \\ c v_n \end{pmatrix}
    \]
\end{definition}

\begin{example}[Scalar Multiplication]
    Let $c = -3$ and $\vct{v} = \begin{pmatrix} 1 \\ -2 \\ 5 \end{pmatrix} \in \R^3$. Then:
    \[
        c\vct{v} = -3 \begin{pmatrix} 1 \\ -2 \\ 5 \end{pmatrix} = \begin{pmatrix} (-3)(1) \\ (-3)(-2) \\ (-3)(5) \end{pmatrix} = \begin{pmatrix} -3 \\ 6 \\ -15 \end{pmatrix}
    \]
\end{example}

\begin{theorem}[Algebraic Properties of $\R^n$]
    For all vectors $\vct{u}, \vct{v}, \vct{w} \in \R^n$ and all scalars $c, d \in \R$, vector addition and scalar multiplication satisfy the following properties:
    \begin{enumerate}[label=\roman*.]
        \item \textbf{Commutativity of Addition:} $\vct{u} + \vct{v} = \vct{v} + \vct{u}$
        \item \textbf{Associativity of Addition:} $(\vct{u} + \vct{v}) + \vct{w} = \vct{u} + (\vct{v} + \vct{w})$
        \item \textbf{Additive Identity:} There exists a zero vector $\vct{0} = \begin{pmatrix} 0 & 0 & \dots & 0 \end{pmatrix}\T \in \R^n$ such that $\vct{v} + \vct{0} = \vct{v}$
        \item \textbf{Additive Inverse:} For each $\vct{v} \in \R^n$, there exists a vector $-\vct{v} = (-1)\vct{v}$ such that $\vct{v} + (-\vct{v}) = \vct{0}$
        \item \textbf{Distributivity over Vector Addition:} $c(\vct{u} + \vct{v}) = c\vct{u} + c\vct{v}$
        \item \textbf{Distributivity over Scalar Addition:} $(c + d)\vct{v} = c\vct{v} + d\vct{v}$
        \item \textbf{Associativity of Scalar Multiplication:} $c(d\vct{v}) = (cd)\vct{v}$
        \item \textbf{Scalar Identity:} $1\vct{v} = \vct{v}$
    \end{enumerate}
\end{theorem}

\subsection{Linear Combinations and Span}

\begin{definition}[Linear Combination]
    Let $\vct{v}_1, \vct{v}_2, \dots, \vct{v}_k \in \R^n$ be vectors and $c_1, c_2, \dots, c_k \in \R$ be scalars. A vector $\vct{w} \in \R^n$ given by
    \[
        \vct{w} = c_1 \vct{v}_1 + c_2 \vct{v}_2 + \dots + c_k \vct{v}_k = \sum_{i=1}^k c_i \vct{v}_i
    \]
    is called a \textbf{linear combination} of the vectors $\vct{v}_1, \vct{v}_2, \dots, \vct{v}_k$ with weights $c_1, c_2, \dots, c_k$.
\end{definition}

% \begin{definition}[Span]
%    The \textbf{span} of a set of vectors $\set{\vct{v}_1, \vct{v}_2, \dots, \vct{v}_k} \subset \R^n$, denoted $\spn\{\vct{v}_1, \dots, \vct{v}_k\}$, is the set of all possible linear combinations of those vectors:
%    \[
%        \spn\{\vct{v}_1, \dots, \vct{v}_k\} = \set{c_1 \vct{v}_1 + c_2 \vct{v}_2 + \dots + c_k \vct{v}_k}{c_1, \dots, c_k \in \R}
%    \]
%\end{definition}

G%eometrically, the span of a single non-zero vector in $\R^n$ represents a line passing through the origin. The span of two non-collinear vectors in $\R^3$ forms a plane passing through the origin.

\begin{exercise}
    Find $c, d \in \R$ such that
    \[
        c \begin{pmatrix} 2 \\ 1 \end{pmatrix} + d \begin{pmatrix} 2 \\ -1 \end{pmatrix} = \begin{pmatrix} 8 \\ 2 \end{pmatrix}
    \]
\end{exercise}

\begin{solution}
    Expanding the linear combination into component equations:
    \begin{align*}
        c \begin{pmatrix} 2 \\ 1 \end{pmatrix} + d \begin{pmatrix} 2 \\ -1 \end{pmatrix} &= \begin{pmatrix} 8 \\ 2 \end{pmatrix} \\ 
        \begin{pmatrix} 2c \\ c \end{pmatrix} + \begin{pmatrix} 2d \\ -d \end{pmatrix} &= \begin{pmatrix} 8 \\ 2 \end{pmatrix} \\ 
        \begin{pmatrix} 2c + 2d \\ c - d \end{pmatrix} &= \begin{pmatrix} 8 \\ 2 \end{pmatrix}
    \end{align*}
    This results in the linear system:
    \begin{align*}
        2c + 2d &= 8 \\ 
        c - d &= 2 
    \end{align*}
    Adding the second equation multiplied by $2$ to the first gives $4c = 12$, so $c = 3$. Substituting $c = 3$ into $c - d = 2$ yields $3 - d = 2 \implies d = 1$.
\end{solution}
\newpage 
%--------------------- 
\begin{exercise}[Describing All Linear Combinations]
    Describe all linear combinations of:
    \begin{enumerate}[label=\arabic*.]
        \item $\vct{a} = \begin{pmatrix} 1 \\ 2 \\ 3 \end{pmatrix}, \quad \vct{b} = \begin{pmatrix} 3 \\ 6 \\ 9 \end{pmatrix}$
        \item $\vct{a} = \begin{pmatrix} 1 \\ 0 \\ 0 \end{pmatrix}, \quad \vct{b} = \begin{pmatrix} 0 \\ 2 \\ 3 \end{pmatrix}$
    \end{enumerate}   
\end{exercise}

\begin{solution}
    \begin{enumerate}[label=\arabic*.]
        \item Since $\vct{b} = 3\vct{a}$, the vectors $\vct{a}$ and $\vct{b}$ are collinear. Any linear combination $c\vct{a} + d\vct{b} = (c + 3d)\vct{a}$ is a scalar multiple of $\vct{a}$, which describes all vectors lying on the line through the origin in direction $\begin{pmatrix} 1 & 2 & 3 \end{pmatrix}\T$.
        \item Since $\vct{a}$ and $\vct{b}$ are not parallel, their linear combinations span a $2$-dimensional plane in $\R^3$ passing through the origin and containing both vectors.
    \end{enumerate}
\end{solution}


\section{Length and Angles from Dot Product}

The dot product introduces metric geometry into Euclidean space $\R^n$, enabling the measurement of lengths, distances, and angles.

\begin{definition}[Dot Product]
    Let $\vct{a}, \vct{b} \in \R^n$. The \textbf{dot product} (or \textbf{Euclidean inner product}) $\vct{a} \cdot \vct{b}$ is the scalar defined by:
    \[
        \vct{a} \cdot \vct{b} \coloneqq \vct{a}\T \vct{b} = \sum_{i=1}^{n} a_i b_i = a_1 b_1 + a_2 b_2 + \dots + a_n b_n \in \R
    \]
\end{definition}

\begin{theorem}[Properties of the Dot Product]
    For all vectors $\vct{a}, \vct{b}, \vct{c} \in \R^n$ and any scalar $k \in \R$:
    \begin{enumerate}[label=\roman*.]
        \item \textbf{Symmetry:} $\vct{a} \cdot \vct{b} = \vct{b} \cdot \vct{a}$
        \item \textbf{Distributivity:} $(\vct{a} + \vct{b}) \cdot \vct{c} = \vct{a} \cdot \vct{c} + \vct{b} \cdot \vct{c}$
        \item \textbf{Homogeneity:} $(k\vct{a}) \cdot \vct{b} = k(\vct{a} \cdot \vct{b})$
        \item \textbf{Positive Definiteness:} $\vct{a} \cdot \vct{a} \ge 0$, with $\vct{a} \cdot \vct{a} = 0$ if and only if $\vct{a} = \vct{0}$.
    \end{enumerate}
\end{theorem}

\subsection{Vector Length and Norm}

\begin{definition}[Length / Norm]
    Let $\vct{a} \in \R^n$. The \textbf{length} (or \textbf{Euclidean norm}) of $\vct{a}$, denoted $\norm{\vct{a}}$, is:
    \[
        \norm{\vct{a}} \coloneqq \sqrt{\vct{a} \cdot \vct{a}} = \sqrt{a_1^2 + a_2^2 + \dots + a_n^2}
    \]
    Note that $\norm{\vct{a}}^2 = \vct{a} \cdot \vct{a}$.
\end{definition}

\begin{definition}[Unit Vector]
    A vector $\vct{u} \in \R^n$ is called a \textbf{unit vector} if $\norm{\vct{u}} = 1$.
    
    Given any non-zero vector $\vct{v} \in \R^n$, dividing by its length yields a unit vector $\vct{u} = \frac{\vct{v}}{\norm{\vct{v}}}$ in the direction of $\vct{v}$ (a process called \textbf{normalization}).
    
    In $\R^2$, any unit vector extending from the origin to the unit circle can be parameterized by an angle $\theta$ as:
    \[
        \vct{u} = \begin{pmatrix} \cos\theta \\ \sin\theta \end{pmatrix}
    \]
\end{definition}

\subsection{Orthogonality and Angles}

\begin{definition}[Orthogonal Vectors]
    Two vectors $\vct{a}, \vct{b} \in \R^n$ are \textbf{orthogonal} (denoted $\vct{a} \perp \vct{b}$) if their dot product is zero:
    \[
        \vct{a} \cdot \vct{b} = 0
    \]
\end{definition}

\begin{remark}[Orthogonal vs. Perpendicular]
    The terms `orthogonal' and `perpendicular' have the same geometric idea. However, `orthogonal' applies to vectors (and includes the zero vector $\vct{0}$, which is orthogonal to every vector), while `perpendicular' is usually for geometric lines and planes.
\end{remark}

\begin{example}[Pythagorean Theorem for Vectors]
    Let $\vct{a}, \vct{b} \in \R^n$ be orthogonal vectors ($\vct{a} \cdot \vct{b} = 0$). Then:
    \[
        \norm{\vct{a} + \vct{b}}^2 = \norm{\vct{a}}^2 + \norm{\vct{b}}^2
    \]
    \begin{proof}
        Expanding the squared norm using dot product properties:
        \begin{align*}
            \norm{\vct{a} + \vct{b}}^2 &= (\vct{a} + \vct{b}) \cdot (\vct{a} + \vct{b}) \\
            &= \vct{a} \cdot \vct{a} + 2(\vct{a} \cdot \vct{b}) + \vct{b} \cdot \vct{b} \\ 
            &= \norm{\vct{a}}^2 + 0 + \norm{\vct{b}}^2 \\ 
            &= \norm{\vct{a}}^2 + \norm{\vct{b}}^2
        \end{align*}
    \end{proof}
\end{example}

\begin{exercise}[Dot Product]
    Let $\vct{a} = \begin{pmatrix} 4 \\ 2 \end{pmatrix}$ and $\vct{b} = \begin{pmatrix} -1 \\ 2 \end{pmatrix}$. Find $\vct{a} \cdot \vct{b}$.
\end{exercise}

\begin{solution}
    Calculating the dot product:
    \begin{align*}
        \vct{a} \cdot \vct{b} &= (4)(-1) + (2)(2) = -4 + 4 = 0
    \end{align*}
    Since $\vct{a} \cdot \vct{b} = 0$, the vectors $\vct{a}$ and $\vct{b}$ are orthogonal.
\end{solution}

\begin{definition}[Angle between Vectors]
    For non-zero vectors $\vct{a}, \vct{b} \in \R^n$, the \textbf{angle} $\alpha \in [0, \pi]$ between $\vct{a}$ and $\vct{b}$ satisfies:
    \[
        \cos\alpha = \frac{\vct{a} \cdot \vct{b}}{\norm{\vct{a}}\norm{\vct{b}}}
    \]
    $\vct{a} \cdot \vct{b} = \norm{\vct{a}}\norm{\vct{b}}\cos\alpha$.
\end{definition}

\begin{theorem}[Vector Inequalities]
    For any vectors $\vct{a}, \vct{b} \in \R^n$:
    \begin{enumerate}[label=\arabic*.]
        \item \textbf{Cauchy-Schwarz Inequality:}
            \[
                \abs{\vct{a} \cdot \vct{b}} \le \norm{\vct{a}}\norm{\vct{b}}
            \]
            Equality holds if and only if $\vct{a}$ and $\vct{b}$ are linearly dependent (collinear).
            
        \item \textbf{Triangle Inequality:}
            \[
                \norm{\vct{a} + \vct{b}} \le \norm{\vct{a}} + \norm{\vct{b}}
            \]
            Equality holds if and only if one vector is a non-negative scalar multiple of the other.
    \end{enumerate}
\end{theorem}

% \begin{definition}[Orthogonal Projection]
 %   Let $\vct{a}, \vct{u} \in \R^n$ with $\vct{u} \neq \vct{0}$. The \textbf{orthogonal projection} of $\vct{a}$ onto $\vct{u}$, denoted $\proj_{\vct{u}}(\vct{a})$, is the vector parallel to $\vct{u}$ given by:
%    \[
%        \proj_{\vct{u}}(\vct{a}) = \left(\frac{\vct{a} \cdot \vct{u}}{\norm{\vct{u}}^2}\right) \vct{u}
%    \]
%\end{definition}


\newpage
\section{Intro to Matrices and Column Spaces}

Matrices provide a structured representation for systems of linear equations, linear transformations, and collections of vectors.

\begin{definition}[Matrix]
    An $m \times n$ \textbf{matrix} $A$ is a rectangular array of real numbers formatted with $m$ rows and $n$ columns:
    \[
        A = \begin{pmatrix}
            a_{11} & a_{12} & \dots & a_{1n} \\
            a_{21} & a_{22} & \dots & a_{2n} \\
            \vdots & \vdots & \ddots & \vdots \\
            a_{m1} & a_{m2} & \dots & a_{mn}
        \end{pmatrix} = [a_{ij}]_{m \times n}
    \]
    Expressed in terms of its column vectors $\vct{a}_1, \vct{a}_2, \dots, \vct{a}_n \in \R^m$, $A$ is written as:
    \[
        A = \begin{pmatrix} \vct{a}_1 & \vct{a}_2 & \dots & \vct{a}_n \end{pmatrix}
    \]
\end{definition}

\subsection{Special Types of Matrices}

Matrices possess structural classifications according to their dimensions and internal symmetry.

\begin{definition}[Matrix Classifications]
    Let $A = [a_{ij}]$ be a matrix.
    \begin{itemize}
        \item \textbf{Square Matrix:} A matrix with equal numbers of rows and columns ($m = n$).
        \item \textbf{Symmetric Matrix:} A square matrix satisfying $A = A\T$, meaning $a_{ij} = a_{ji}$ for all $i, j$.
        \item \textbf{Diagonal Matrix:} A square matrix where all off-diagonal entries are zero ($a_{ij} = 0$ for all $i \neq j$).
        \item \textbf{Identity Matrix ($I_n$):} A diagonal matrix with ones on the main diagonal ($a_{ii} = 1$) and zeros elsewhere.
        \item \textbf{Triangular Matrices:}
            \begin{itemize}
                \item \textbf{Upper Triangular Matrix:} A square matrix with zeros below the main diagonal ($a_{ij} = 0$ for $i > j$).
                \item \textbf{Lower Triangular Matrix:} A square matrix with zeros above the main diagonal ($a_{ij} = 0$ for $i < j$).
            \end{itemize}
    \end{itemize}
\end{definition}

\newpage
\begin{example}[Types of Matrices]
    \begin{itemize}
        \item \textbf{3$\times$3 Identity Matrix (Symmetric and Diagonal):}
            \[
                I_3 = \begin{pmatrix} 1 & 0 & 0 \\ 0 & 1 & 0 \\ 0 & 0 & 1 \end{pmatrix}
            \]
        \item \textbf{3$\times$3 Diagonal Matrix:}
            \[
                D = \begin{pmatrix} 2 & 0 & 0 \\ 0 & 4 & 0 \\ 0 & 0 & 5 \end{pmatrix}
            \]
        \item \textbf{Triangular Matrices:}
            \begin{itemize}
                \item 2$\times$2 Upper Triangular Matrix:
                    \[
                        U = \begin{pmatrix} 1 & 2 \\ 0 & 3 \end{pmatrix}
                    \]
                \item 3$\times$3 Lower Triangular Matrix:
                    \[
                        L = \begin{pmatrix} 1 & 0 & 0 \\ 2 & 3 & 0 \\ 3 & 1 & 5 \end{pmatrix}
                    \]
            \end{itemize}
    \end{itemize}
\end{example}

\subsection{Matrix-Vector Product}

\begin{definition}[Matrix-Vector Product]
    Let $A = [a_{ij}]$ be an $m \times n$ matrix with rows $\vct{r}_1\T, \vct{r}_2\T, \dots, \vct{r}_m\T \in \R^n$ and column vectors $\vct{a}_1, \vct{a}_2, \dots, \vct{a}_n \in \R^m$. For a vector $\vct{x} = \begin{pmatrix} x_1 & x_2 & \dots & x_n \end{pmatrix}\T \in \R^n$, the \textbf{matrix-vector product} $A\vct{x} \in \R^m$ is defined componentwise by:
    \[
        A\vct{x} = \begin{pmatrix} \vct{r}_1 \cdot \vct{x} \\ \vct{r}_2 \cdot \vct{x} \\ \vdots \\ \vct{r}_m \cdot \vct{x} \end{pmatrix}
    \]
\end{definition}

\begin{remark}[Linear Combination Viewpoint]
    Equivalently, the product $A\vct{x}$ is a \textbf{linear combination of the column vectors of $A$} with weights given by the entries of $\vct{x}$:
    \[
        A\vct{x} = x_1 \vct{a}_1 + x_2 \vct{a}_2 + \dots + x_n \vct{a}_n
    \]
    This dual perspective links matrix multiplication directly to linear combinations and vector spans.
\end{remark}

\begin{exercise}[Matrix-Vector Product]
    Let $A = \begin{pmatrix} 1 & 3 \\ 2 & 4 \end{pmatrix}$ and $\vct{h} = \begin{pmatrix} 1 \\ 2 \end{pmatrix}$. Find $A\vct{h}$.
\end{exercise}

\begin{solution}
    Using the row dot-product formula:
    \begin{align*}
        A\vct{h} &= \begin{pmatrix} \begin{pmatrix} 1 & 3 \end{pmatrix} \cdot \begin{pmatrix} 1 & 2 \end{pmatrix}\T \\[4pt] \begin{pmatrix} 2 & 4 \end{pmatrix} \cdot \begin{pmatrix} 1 & 2 \end{pmatrix}\T \end{pmatrix} \\
        &= \begin{pmatrix} (1)(1) + (3)(2) \\ (2)(1) + (4)(2) \end{pmatrix} = \begin{pmatrix} 7 \\ 10 \end{pmatrix}
    \end{align*}
    Alternatively, using column linear combinations:
    \[
        A\vct{h} = 1\begin{pmatrix} 1 \\ 2 \end{pmatrix} + 2\begin{pmatrix} 3 \\ 4 \end{pmatrix} = \begin{pmatrix} 1 \\ 2 \end{pmatrix} + \begin{pmatrix} 6 \\ 8 \end{pmatrix} = \begin{pmatrix} 7 \\ 10 \end{pmatrix}
    \]
\end{solution}

\subsection{Column Dependence} 
\begin{example}
    \[
      A = \begin{pmatrix}
        1 & 2 & 3 \\ 1 & 4 & 5 \\ 6 & 0 & 6
      \end{pmatrix}
    \]
    Are columns dependent? 

    Let's check if a linear combination of column 1 and 2 can result in 3 

    We can see $\text{col}1 + \text{col}2 = \text{col}3$

    So column 3 is dependent. 

    Another way of writing this is asking:

    If we create a system of equations with the other columns, can it result in this one? 

    If the system of equations has a solution, the column is independent. 
\end{example}

\newpage
\subsection{Column Space}
\begin{definition}[Column Space]
    The column space of $A$ ($C(A)$) is defined as: 

    \[
        C(A) = \{A\vct{x} \mid \vct{x} \in \R^n\}
    \]

    Or: 
    \[
        A = (\vct{a}_1, \hdots \vct{a}_n) 
    \]
    \[
        C(A) = \{x_1 \vct{a}_1 + \hdots + x_n \vct{a}_n \mid x_1 \hdots x_n \in \R\}
    \]
\end{definition}

This explanation is kind of confusing. The column space is every possible destination you can reach by scaling and combining the columns of the matrix. So, $C(A)$ is the set of every possible destination in this matrix. A single non-zero column is a 1D line passing through the origin. Two independent columns is a flat 2D plane passing through the origin, and 3 independent columns is a solid 3D volume.

\begin{example}
    \[
        A_1 = \begin{pmatrix}
          1 & 0 & 0 \\ 0 & 1 & 0 \\ 0 & 0 & 1 
        \end{pmatrix}
    \]
    What space does $A_1$ span? 

    \[
        C(A_1) = \{A_1\vct{x} \mid \vct{x} \in \R^3\}
    \]

    $C(A_1)$ spans $\R^3$ because there are 3 independent columns.
    
\end{example}

\begin{example}
    \[
        A_2 = \begin{pmatrix}
            1 & 0 & 0 \\ 0 & 1 & 0 \\ 0 & 0 & 0
        \end{pmatrix}
    \]
    What does $A_2$ span? 
    
    $A_2$ spans a plane in $\R^3$ containing the points $(1, 0, 0)$ and $(0, 1, 0)$
    
\end{example}


\subsubsection{Rank}
\begin{definition}
    Let $A$ be an $m \times n$ matrix. The rank of $A$ is equal to the number of independent columns $A$ has.

    For example, $\text{Rank}(A_1) = 3$ and $\text{Rank}(A_2) = 2$ (From the earlier example). 

    It is also denoted as $\text{rk}(A)$
\end{definition}


\newpage

\section{Matrix-Matrix Product}

\begin{definition}
    Let $A$ be an $m \times r$ matrix and $B$ be an $r \times n$ matrix. 
    
    The Matrix-Matrix product $AB$ is an $m \times n$ matrix. 
    
    $B = (\vct{b}_1 \hdots \vct{b}_n)$

    $AB = (A\vct{b}_1 \hdots A\vct{b}_n)$
\end{definition}

\begin{example}
    \begin{gather*}
        A = \begin{pmatrix}
            1 & 2 \\ 3 & 4
        \end{pmatrix}, \,  
        B = \begin{pmatrix}
            0 & 1 \\ 1 & 0 
        \end{pmatrix} = \begin{pmatrix}
            \vct{b}_1 & \vct{b}_2
        \end{pmatrix}
         \\ 
        A\vct{b}_1 = \begin{pmatrix}
            1 & 2 \\ 3 & 4 
        \end{pmatrix} \begin{pmatrix}
          0 \\ 1 
        \end{pmatrix} \\ 
        A\vct{b}_1 = \begin{pmatrix}
            2 \\ 4 
        \end{pmatrix} \\ 
        A\vct{b}_2 = \begin{pmatrix}
            1 & 2 \\ 3 & 4 
        \end{pmatrix} \begin{pmatrix}
            1 \\ 0
        \end{pmatrix} \\ 
        A\vct{b}_2 = \begin{pmatrix}
            1 \\ 3 
        \end{pmatrix} \\ 
        AB = \begin{pmatrix}
            2 & 1 \\ 4 & 3 
        \end{pmatrix}
    \end{gather*}
\end{example}

\subsubsection{Properties of Matrix-Matrix multiplication}

Matrices are an algebra, just like natural numbers.
However, they lose commutativity.
\begin{itemize}
  \item Non-commutative
      $AB \neq BA$ in most cases. 
  \item Associative 
      $(AB)C = A(BC)$
  \item Distributive 
    $A(B+C) = AB+AC$ and $(B+C)A = BA+CA$
\end{itemize}

\newpage
\subsubsection{Rank 1 Matrices}
Matrices with rank 1 are called Rank 1 Matrices. They can be described as the product of a column and row matrix. 


\begin{example}
    
    \[
    \text{Let } A = \begin{pmatrix}
        1 & 2 & 10 & 100 \\ 3 & 6 & 30 & 300 \\ 2 & 4 & 20 & 200
    \end{pmatrix}
    \]

    Since $A$ is Rank 1, we can write it as the product of a column and a row matrix. 

    \[
        A = \begin{pmatrix}
          1 \\ 3 \\ 2 
        \end{pmatrix} \begin{pmatrix}
        1 & 2 & 10 & 100
        \end{pmatrix}
    \]
\end{example}

This is also called $CR$ decomposition.

\begin{definition}[CR Decomposition]
    $A$ is an $m \times n$ matrix.

    $C$ is an $m \times 1$ matrix containing a non-zero column of $A$

    $R$ is a $1 \times n$ matrix containing the multiples/weights of the columns in $C$

    $A = CR$ 
\end{definition}

\begin{example}
    The previous example can also be written as:
    \[
        A = \begin{pmatrix}
            2 \\ 6 \\ 4
        \end{pmatrix} \begin{pmatrix}
        \frac{1}{2} & 1 & 5 & 50
        \end{pmatrix}
    \]
    
\end{example}

